\chapter{全同粒子、对称性与守恒律}
\section{全同粒子}
上一节，我们讨论了粒子的典型系统，接下来，我们讨论多粒子的情况. 考虑两个不同粒子，它们的态矢量分别为\(\ket{\al},\ket{\be}\)，假设二者可分辨，我们可以用\(\ket{\al;\be}\)
表征. 然而，真实世界的粒子实际上是“不可分辨的”的，我们不能把电子或其他粒子想象成一个小球，因此我们也不能用经典的手段去区分一个粒子和另外一个，这就导致两个粒子的交换
往往是不可察觉的. 如果强行区分两个粒子，那就可能导致同一个能量下出现两种不同的状态，即12和21是不同的，这也就导致了简并，这种简并被称为\textbf{交换简并}. 我们不妨定义
交换算符\(P_{12}\)，这表示我们把人为编号1的粒子和人为编号为2的粒子交换位置，当交换算符作用于两个粒子组合的右矢，会得到：
\[
P_{12}\ket{\al;\be} = \ket{\be;\al}
\]
利用这个性质，我们很容易得到：
\begin{align*}
& P_{12} = P_{21} \\
& P_{12}^{2} = 1 
\end{align*}
由于假定可区分的粒子的右矢具备与无耦合表象的角动量本征矢类似的性质，这时我们可以定义对1的可观测量和对2的可观测量分别为：
\[
A_{1}\otimes I_{2}\ , \ I_{1}\otimes A_{2}
\]
让可观测量对态矢作用，有：
\begin{align*}
& A_{1}\ket{a_{1};a_{2}} = a_{1}\ket{a_{1};a_{2}} \\
& A_{2}\ket{a_{1};a_{2}} = a_{2}\ket{a_{1};a_{2}}
\end{align*}
现在两边左乘\(P_{12}\)，有：
\begin{align*}
& P_{12}A_{1}P_{12}^{-1}P_{12}\ket{a_{1};a_{2}} = a_{1}P_{12}\ket{a_{1};a_{2}} \\
\therefore\ 
& P_{12}A_{1}P_{12}^{-1}\ket{a_{2};a_{1}} = a_{1}\ket{a_{2};a_{1}} \\
\end{align*}
这样我们得到了交换算符对可观测量的影响：
\[
A_{2} = P_{12}A_{1}P_{12}^{-1}
\]
那么对于全同粒子的Hamiltonian：
\[
H = \frac{\vp_{1}^{2}}{2m} + \frac{\vp_{2}^{2}}{2m} + U_{in}(\abs{\vr_{1}-\vr_{2}}) + U_{ext}(\vr_{1}) + U_{ext}(\vr_{2})
\]
我们做交换后可以得到：
\[
H = P_{12}HP_{12}^{-1} \Longleftrightarrow P_{12}H = HP_{12}
\]
这也就意味着交换算符是和Hamiltonian对易的，根据Heisenberg方程，交换算符是守恒的，可以将其看作运动常数. 由\(P_{12}^{2} = 1\)可得，交换算符的本征值为\(\pm1\). 这时，
我们可以构造出交换算符的本征态：
\begin{align*}
& \ket{a_{1}a_{2}}_{+} = \frac{1}{\sqrt{2}}\br{\ket{a_{1};a_{2}} + \ket{a_{2};a_{1}}} \\
& \ket{a_{1}a_{2}}_{-} = \frac{1}{\sqrt{2}}\br{\ket{a_{1};a_{2}} - \ket{a_{2};a_{1}}}
\end{align*}
这个时候，交换算符的本征态分别对应了一个对称的和一个反对称的，这两种本征态分别对应了两种不同的粒子：前者对应Boson子，后者为Fermi子. 推广到N个粒子也是这样. Boson子和
Fermi子的性质是有区别的，对于前者，两个粒子不可能处于同一个状态，因为如果\(\al=\beta\)的话，反对称的本征右矢实际上等于\(0\)，这就意味着不可能存在两个全同Fermi子处于相同
状态. 这种规则被称为\textbf{Pauli不相容原理}. 对于Boson子，它可以存在3中可能的态：
\[
\ket{\al;\al}\ , \ \ket{\be;\be}\ , \ \frac{1}{\sqrt{2}}\br{\ket{\al;\be}+\ket{\be;\al}}
\]
同时，在非相对论量子力学的语境下，Fermi子对应的是半整数自旋，而Boson子对应的是整数自旋. 明确了两种粒子态的基本性质，我们来看波函数的特点. 一个粒子的状态不仅由位置确定，
还由其自旋确定，所以对于双粒子的情况，其坐标表象下的波函数可以用\(\bra{\vr,m_{s}}\)得到：
\[
\bra{\vr_{1},m_{s1};\vr_{2},m_{s2}}\ket{\al;\be} = \sum\bra{\vr_{1},m_{s1};\vr_{2},m_{s2}}\ket{a_{1};a_{2}}\bra{a_{1};a_{2}}\ket{\al,\be}
\]
从函数的视角来看这个问题，我们实际上将代表这两个粒子状态的波函数分成了两个组成部分，空间部分和自旋部分：
\[
\psi_{12} = \phi(\vr_{1},\vr_{2})\chi_{12}
\]
其中\(\chi_{12}\)是两个电子耦合后的自旋角动量. 当两个电子平行时，叠加后产生三重态；当两个电子反平行时，叠加后是单态. 对于前者，叠加后形式对称；后者是反对称的. 这时做交换后
波函数中自旋满足：
\[
\chi_{21} = (-1)^{S-1}\chi_{12}
\]
同时，空间部分的波函数的自变量\(\vr_{1},\vr_{2}\)改写为\(\vr_{C},\vr=\vr_{1}-\vr_{2}\). 这样我们的空间部分波函数实际上可以改写为：
\[
\phi(\vr_{1},\vr_{2}) = \phi(\vr) = \Rcal(r)Y_{L,m_{L}}(\uv{n})
\]
利用交换算符作用，作用后，有：
\[
Y_{L,m_{L}}(-\uv{n}) = (-1)^{L}Y_{L,m_{L}}(\uv{n})
\]
于是，电子交换后的波函数满足：
\[
P_{12}\psi_{12} = (-1)^{L+S-1}\psi_{12}
\]
对于空间的波函数，我们还有一种构造手法是：
\[
\phi(\vr_{1},\vr_{2}) = \psi_{1}(\vr_{1})\psi_{2}(\vr_{2}) \pm \psi_{1}(\vr_{2})\psi_{2}(\vr_{1})
\]
考虑1维情况，计算这样一个均值，有：
\begin{align*}
\avg{(x_{1}-x_{2})^{2}}
& = \int\dd x_{1}\int\dd x_{2}(x_{1} - x_{2})^{2}\abs{\phi(x_{1},x_{2})}^{2} \\
& = \avg{x_{1}^{2}} + \avg{x_{2}^{2}}-2\avg{x_{1}}\avg{x_{2}}\pm2\int\dd{x}_{1}\int\dd{x}_{2}\br{x_{1}\psi_{1}^{*}\psi_{2}}\br{x_{2}\psi_{1}\psi^{*}_{2}}
\end{align*}
我们会发现Boson子的平均距离会靠的更近，而Fermi子恰好相反，这种表现仿佛是Boson子内部和Fermi子内部存在某种相互作用. 这种相互作用被称为\textbf{交换力}. 交换力纯粹是
量子力学效应，没有经典对应. 